\documentclass[11pt]{article}
\newcommand{\SheetCode}{Theory 09}
\usepackage[margin=25mm]{geometry}
\usepackage[T1]{fontenc}
\usepackage{lmodern}
\DeclareUnicodeCharacter{2605}{$\star$}
\usepackage{microtype}
\usepackage{amsmath,amssymb,mathtools,bm}
\usepackage{enumitem}
\usepackage{xcolor}
\usepackage{fancyhdr}
\usepackage{hyperref}
\usepackage[most]{tcolorbox}
\definecolor{agiBlue}{HTML}{173B57}
\definecolor{agiGold}{HTML}{B7791F}
\definecolor{agiLight}{HTML}{EFF6FA}
\definecolor{agiGreen}{HTML}{184D3B}
\hypersetup{colorlinks=true,linkcolor=agiBlue,urlcolor=agiBlue}
\setlength{\parindent}{0pt}
\setlength{\parskip}{0.55em}
\setlength{\headheight}{14pt}
\setlist{nosep,leftmargin=*}
\pagestyle{fancy}
\fancyhf{}
\fancyhead[L]{\small Part of the \href{https://github.com/mmikhasenko/GIModel.jl}{Agentic Gottried--Isgur} project}
\fancyhead[R]{\SheetCode}
\fancyfoot[C]{\thepage}
\newcommand{\dd}{\,\mathrm d}
\newcommand{\ket}[1]{\lvert #1\rangle}
\newcommand{\bra}[1]{\langle #1\rvert}
\newcommand{\braket}[2]{\langle #1\mid #2\rangle}
\newcommand{\expect}[1]{\left\langle #1\right\rangle}
\newcommand{\op}[1]{\widehat{#1}}
\newcommand{\GeV}{\,\mathrm{GeV}}
\newcommand{\R}{\mathbb{R}}
\newtcolorbox{goals}{colback=agiLight,colframe=agiBlue,title=Learning outcomes}
\newtcolorbox{reading}{colback=white,colframe=agiGold,title=Book reference,
  after skip=5mm}
\newtcolorbox{paperbridge}{colback=agiGreen!7,colframe=agiGreen,
  colbacktitle=agiGreen,coltitle=white,title=Connection to the GI paper}
\newtcolorbox{problem}[1]{colback=white,colframe=agiBlue,title={#1},
  before skip=3.5mm,after skip=1mm}
\newtcolorbox{solution}{colback=black!2,colframe=black!45,title=Solution}
\newtcolorbox{quiz}{colback=white,colframe=agiGold,title=Post-solution concept check}
\newtcolorbox{checkpoint}{colback=white,colframe=agiGold,title=Interpretation checkpoint}
\newcommand{\sheettitle}[3]{%
  \begin{center}
  {\Large\bfseries #1}\par
  \vspace{0.2em}{\color{agiBlue}#2}\par
  \vspace{0.4em}\small #3
  \end{center}\vspace{0.5em}}
\newcommand{\difficulty}[1]{\textbf{Difficulty:} #1}
\newcommand{\timebox}[1]{\textbf{Suggested time:} #1}
\newcommand{\answer}{\par\smallskip\color{black!50}\textbf{Answer:} }

\begin{document}
\sheettitle{Sheet 9: Reconstruct the GI computation}{A paper-order capstone and audit}{Prerequisite: Sheets 1--8. \difficulty{★★★} \quad \timebox{3--4 hours}}

\begin{goals}
Assemble the entire calculation in dependency order; distinguish physics from
representation; design convergence and invariant tests; diagnose disagreements;
explain what it means to reproduce a published spectrum without retuning.
\end{goals}
\begin{reading}
For a selective review in Landau--Lifshitz, \emph{Quantum Mechanics:
Non-Relativistic Theory}, use Chs. I--VI for operators, the Schr\"odinger
equation, angular momentum, central fields, and perturbation theory; add
Chs. VIII and XIV for spin and angular-momentum coupling. This is a reference
map, not a request to reread or follow the complete book.
\end{reading}


\begin{problem}{Problem 1 ★★★ — Draw the complete dependency graph}
Starting from $(m_q,b,c,a_k,\gamma_k,\sigma_0,s,\epsilon_i)$ and a requested
$n^{2S+1}L_J$ list, order these operations: central smearing, running coupling,
fixed-channel diagonalization, momentum sandwiches, spectroscopic mixing,
flavor annihilation, decay quantities, and reference comparison. Justify every
arrow. A dependency graph is a diagram in which an arrow $A\to B$ means that
$A$ must be known before $B$ can be constructed.

\textbf{Hint:} Use the definitions in Sheets 3--8; $a_k$ and $\alpha_k$
name the same coupling coefficients. This is a symbolic dependency exercise,
not a numerical fit. Include the annihilation strengths of Sheet 7 and the
chosen decay operators, charges/couplings, and final-state masses as additional
inputs where needed. Label the outputs of each stage and the separate source
of published comparison data.
\end{problem}
\begin{solution}
Parameters define $\alpha_s$, mass-dependent $\sigma$, smeared kernels, and
momentum factors. Together with kinematics and angular coefficients they define
each fixed-$(L,S,J)$ Hamiltonian. Diagonalization produces masses and radial
waves in the representation used by each numerical method. Those waves define
off-diagonal tensor/antisymmetric spin--orbit blocks;
their diagonalization produces spectroscopic components. Self-conjugate
channels then acquire flavor-annihilation blocks. Physical component waves feed
decay quantities. Only after this ordered physics calculation should state names and
published rows enter the comparison layer. Reversing an arrow either uses a
wave before it exists or contaminates the model with reference data.
\end{solution}
\begin{quiz}
Where is the boundary between prediction and reproduction audit? \answer the
model ends at physical masses/waves/observables; digitized paper rows and
residuals belong to a separate comparison layer.
\end{quiz}

\begin{problem}{Problem 2 ★★★ — Prove a representation contract}
Let $U$ be a unitary map from an HO coefficient representation to another
complete representation. Show how $\ket\psi$, $H$, and $O$ transform and prove
that eigenvalues and expectations are invariant.

\textbf{Hint:} Use $U^\dagger U=UU^\dagger=I$ and
$\ket{\psi'}=U\ket\psi$. Require the transformed operators to act on the
transformed domain. An arbitrary pair of finite grid and HO truncations need
not be related by an exact unitary map; this proof concerns equivalent complete
representations.
\end{problem}
\begin{solution}
$\ket{\psi'}=U\ket\psi$, $H'=UHU^\dagger$, and $O'=UOU^\dagger$. Then
$H'\ket{\psi'}=EU\ket\psi=E\ket{\psi'}$ and
$\bra{\psi'}O'\ket{\psi'}=\bra\psi O\ket\psi$. Stored arrays need not match;
the invariant operations are norms, overlaps, expectations, and spectra.
\end{solution}
\begin{quiz}
What must a new numerical method for the radial equation provide before it can
replace an existing method? \answer the same physical operations needed later
--- norms, overlaps, expectation values, and momentum-space operations --- not
the internal array layout used by one existing method.
\end{quiz}

\begin{problem}{Problem 3 ★★★ — Design a layered validation suite}
Using the identities and examples in Sheets 1--8, propose one test for each
layer: radial interactions, angular algebra, operator assembly, eigenvalue
solution, mixing, and comparison with published results.
Present a six-row table with columns: input or toy case, quantity checked,
expected result, and likely error detected. Give exact identities where
available and state which numerical control is varied for convergence tests.
A residual is a calculated value minus its published reference.
Explain why one overall root-mean-square residual cannot replace these tests.

\textbf{Hint:} Start with an analytic limit or invariant for each layer.
\end{problem}
\begin{solution}
One possible suite is below. Numerical tolerances must reflect the units and
accuracy target of each quantity.

{\small
\begin{tabular}{@{}p{0.20\linewidth}p{0.25\linewidth}p{0.23\linewidth}p{0.23\linewidth}@{}}
\textbf{Input / case} & \textbf{Quantity checked} & \textbf{Expected result} & \textbf{Error detected} \\[3pt]
Gaussian kernel & $\int\rho_\sigma\dd^3r$ & $1$ & Missing prefactor or volume factor \\[5pt]
Common-wave $^3P_J$ & Weighted spin--orbit sum & $0$ & Wrong angular or degeneracy factor \\[5pt]
Hermitian $B,V$ & $(BVB)^\dagger-BVB$ & Zero matrix & Wrong operator order or adjoint \\[5pt]
Fixed physical $H$ & Energies under $N,\beta$ refinement & Stable within $\delta E$ & Basis truncation \\[5pt]
Real $2\times2$ block & Eigenvalue sum and product & $\operatorname{tr}M$, $\det M$ & Incorrect diagonalization \\[5pt]
Published mass rows & Signed residual by state and flavor & Matches declared reproduction tolerance & Wrong inputs or row assignment
\end{tabular}}

A single RMS value can hide a failing sector among many accurate states,
and cannot identify which layer caused the discrepancy. Analytic and
convergence checks remain necessary even when the final mass residuals are small.
\end{solution}
\begin{quiz}
Which test most directly detects a stray $4\pi$ in radial normalization?
\answer an independently normalized three-dimensional kernel combined with a
known normalized reduced radial wave and a scale-invariance check.
\end{quiz}

\begin{problem}{Problem 4 ★★★ — Diagnose three disagreements}
All physical model parameters are held fixed. For each symptom below,
propose one initial hypothesis, one controlled check, and the outcome that
would support or weaken your hypothesis:
\begin{enumerate}[label=(\alph*)]
\item All central masses change when the grid spacing is reduced at fixed box size.
\item The $^3P_J$ centroid agrees with a reference but the splittings are smaller.
\item Nonmixing masses agree with a reference while isoscalar masses disagree.
\end{enumerate}
These symptoms need not have a unique cause; justify your first check.

\textbf{Hint:} Follow the dependencies in Problem 1 and vary one ingredient at a time.
\end{problem}
\begin{solution}
(a) Suspect the derivative or kinetic discretization. Halve the spacing at
fixed box size and compare with an independently converged HO result.
Systematic approach to a common energy supports a discretization explanation;
a persistent discrepancy calls for checks of operator assembly and box size.

(b) Suspect the tensor or spin--orbit contributions. First compare their
angular coefficients with Sheet 6, then compare their radial matrix elements
using one common normalized wave. Incorrect angular values support an algebra
error; correct angular values with smaller radial contributions point to the
kernels or momentum factors. A stable centroid alone cannot decide between them.

(c) Suspect the flavor-annihilation block or reference assignment.
Hold the converged nonmixing states fixed and compare the block entries with
the Sheet-7 construction, including flavor factors and wave functionals.
A mismatch localizes the problem to this stage; agreement shifts attention
to the physical inputs and state-to-reference matching.
\end{solution}
\begin{quiz}
Why is retuning the first response scientifically weak? \answer it can absorb
an implementation error and destroys the claim of reproducing published inputs.
\end{quiz}

\begin{problem}{Problem 5 ★★★ — Oral defense: explain GI in ten sentences}
Prepare ten sentences, each with one of these verbs: define, reduce, classify,
smear, relativize, diagonalize, mix, observe, validate, reproduce. Your answer
must connect every verb to a mathematical object.
\end{problem}
\begin{solution}
One model answer: Define constituent masses and one common parameter set. Reduce
the rest-frame two-body problem to relative radial and angular sectors. Classify
states by $n^{2S+1}L_J$ and exact discrete symmetries. Smear singular kernels
with a normalized mass-dependent Gaussian. Relativize inverse-mass interactions
with Hermitian functions of $p^2$ on both sides. Diagonalize the complete
fixed-$(L,S,J)$ Hamiltonian in a converged HO basis (and compare with an FD
representation). Mix converged channels through tensor and antisymmetric
spin--orbit blocks, then mix self-conjugate flavors through annihilation.
Observe masses and decay functionals using the resulting component waves.
Validate analytic identities, normalization, Hermiticity, convergence, and
cross-representation invariants. Reproduce the paper only by comparing these
predictions to independently curated reference data without retuning.
\end{solution}
\begin{quiz}
Final test: which single word prevents the course from becoming a list of
formulas? \answer dependency. Understanding GI means knowing why each object is
needed before the next one can be defined.
\end{quiz}

\begin{paperbridge}
GI Eq. (14) and its following paragraphs state the three-stage spectrum
algorithm: fixed-$\ket{jm;ls}$ diagonalization, tensor/antisymmetric spin--orbit
mass matrices, and isoscalar annihilation. Table II supplies the common model
parameters; Figs. 3--11 display the predicted sectors; Sec. IV and Tables IV--VII
test the resulting wavefunctions through couplings and decay observables.
\end{paperbridge}
\end{document}
